\documentclass[tikz,border=3pt]{standalone}

\usepackage[UTF8]{ctex}
\usepackage{tikz}
\usepackage{xcolor}
\usetikzlibrary{positioning,arrows.meta,calc,fit,backgrounds}

% ===== Colors: aligned with existing thesis figures =====
\definecolor{capblue}{RGB}{222,237,252}
\definecolor{slotcyan}{RGB}{218,244,241}
\definecolor{mdborange}{RGB}{253,235,211}
\definecolor{cdtyellow}{RGB}{255,246,198}
\definecolor{mgrgreen}{RGB}{225,242,218}
\definecolor{opred}{RGB}{248,222,222}
\definecolor{framegray}{RGB}{70,70,70}
\definecolor{textgray}{RGB}{35,35,35}

\begin{document}

\begin{tikzpicture}[
    font=\small,
    text=textgray,
    >=Latex,
    group/.style={
        rectangle,
        rounded corners=5pt,
        draw=framegray,
        line width=0.9pt,
        inner sep=5pt
    },
    groupTitle/.style={
        anchor=west,
        align=left,
        text=framegray
    },
    cell/.style={
        rectangle,
        rounded corners=3pt,
        draw=framegray!80,
        line width=0.65pt,
        fill=white,
        align=center,
        minimum height=11mm,
        inner xsep=3pt,
        inner ysep=2pt
    },
    maincell/.style={
        cell,
        text width=122mm
    },
    subcell/.style={
        cell,
        text width=37mm,
        minimum height=14mm
    },
    sidecell/.style={
        cell,
        text width=43mm,
        minimum height=15mm
    },
    arrow/.style={
        -{Latex[length=2.2mm,width=1.7mm]},
        draw=framegray,
        line width=0.85pt
    },
    softarrow/.style={
        -{Latex[length=2.0mm,width=1.5mm]},
        draw=framegray!72,
        line width=0.7pt
    },
    note/.style={
        font=\footnotesize,
        align=center,
        text=framegray
    }
]

% ===== Main layered mapping =====
\node[maincell] (cap) at (0, 6.45) {
    \textbf{能力语义层}\quad
    \texttt{capability}\quad
    \footnotesize 定义能力类型、对象、权限、badge、untyped 与 zombie 语义，是 CSpace 证明的能力语义基底
};

\node[maincell, below=10mm of cap] (cte) {
    \textbf{槽位对象层}\quad
    \texttt{cspace::cte}\quad
    \footnotesize 将 \texttt{cap + mdb\_node} 组织为槽位对象，承接局部规格、最终能力判断与桥接义务
};

\node[subcell] (mdb) at (-4.45, 1.10) {
    \textbf{MDB 图层}\\[-0.4mm]
    \texttt{cspace::mdb}\\[-0.3mm]
    \footnotesize 维护活跃槽位与\\
    \footnotesize 前驱后继图编辑
};

\node[subcell] (cdt) at (4.45, 1.10) {
    \textbf{CDT 派生层}\\[-0.4mm]
    \texttt{cspace::cdt}\\[-0.3mm]
    \footnotesize 显式表达 parent\_of、\\
    \footnotesize is\_original 与派生树义务
};

% ===== Manager core group =====
\node[subcell] (spec) at (-4.85, -2.45) {
    \textbf{规格层}\\[-0.4mm]
    \texttt{manager::spec}\\[-0.3mm]
    \footnotesize 定义总 \texttt{wf}、\\
    \footnotesize 跨层关系与删除规格
};

\node[subcell] (impl) at (0, -2.45) {
    \textbf{执行层}\\[-0.4mm]
    \texttt{manager::impl\_*}\\[-0.3mm]
    \footnotesize 承载 insert / move / swap /\\
    \footnotesize delete / revoke 核心路径
};

\node[subcell] (proof) at (4.85, -2.45) {
    \textbf{证明层}\\[-0.4mm]
    \texttt{manager::proof}\\[-0.3mm]
    \footnotesize 分阶段恢复 MDB / CDT，\\
    \footnotesize 组合得到管理器层 \texttt{wf}
};

% ===== Side connection layers =====
\node[sidecell] (resolve) at (-4.95, -6.25) {
    \textbf{解析层}\\[-0.4mm]
    \texttt{cspace::resolve}\\[-0.3mm]
    \footnotesize 给出 \texttt{resolve\_address\_bits}\\
    \footnotesize 及故障分支的抽象解析语义
};

\node[sidecell] (kernel) at (4.95, -6.25) {
    \textbf{内核连接层}\\[-0.4mm]
    \texttt{cspace::kernel}\\
    \footnotesize 与兼容封装层一起连接\\
    \footnotesize kernel 入口与管理器实例化
};

\node[note] (boundary) at (0, -8.45) {
    \textbf{主结论边界}\\
    \footnotesize 最强机械化结论集中在 \texttt{CSpaceManager::*} 核心层，\\
    \footnotesize 解析层形成独立查找路线，kernel 连接层暂不作为最强结论层
};

% ===== Title padding anchors =====
\coordinate (capPad) at ($(cap.north)+(0,6.4mm)$);
\coordinate (ctePad) at ($(cte.north)+(0,5.6mm)$);
\coordinate (mdbPad) at ($(mdb.north)+(0,6.4mm)$);
\coordinate (cdtPad) at ($(cdt.north)+(0,6.4mm)$);
\coordinate (mgrPad) at ($(spec.north)+(0,6.4mm)$);
\coordinate (resolvePad) at ($(resolve.north)+(0,6.4mm)$);
\coordinate (kernelPad) at ($(kernel.north)+(0,6.4mm)$);
% ===== Background groups =====
\begin{scope}[on background layer]
\node[group, fill=capblue, fit=(capPad)(cap)] (gcap) {};
\node[group, fill=slotcyan, fit=(ctePad)(cte)] (gcte) {};
\node[group, fill=mdborange, fit=(mdbPad)(mdb)] (gmdb) {};
\node[group, fill=cdtyellow, fit=(cdtPad)(cdt)] (gcdt) {};
\node[group, fill=mgrgreen, fit=(mgrPad)(spec)(impl)(proof)] (gmgr) {};
\node[group, fill=opred, fit=(resolvePad)(resolve)] (gresolve) {};
\node[group, fill=opred, fit=(kernelPad)(kernel)] (gkernel) {};
\end{scope}

% ===== Group titles =====
\node[groupTitle] at ($(gcap.north west)+(2mm,-2.5mm)$) {
    \textbf{基础语义层}
};
\node[groupTitle] at ($(gcte.north west)+(2mm,-2.5mm)$) {
    \textbf{槽位抽象层}
};
\node[groupTitle] at ($(gmdb.north west)+(2mm,-2.5mm)$) {
    \textbf{结构层}
};
\node[groupTitle] at ($(gcdt.north west)+(2mm,-2.5mm)$) {
    \textbf{派生层}
};
\node[groupTitle] at ($(gmgr.north west)+(2mm,-2.5mm)$) {
    \textbf{管理器核心层}
};
\node[groupTitle] at ($(gresolve.north west)+(2mm,-2.5mm)$) {
    \textbf{解析路线}
};
\node[groupTitle] at ($(gkernel.north west)+(2mm,-2.5mm)$) {
    \textbf{内核桥接}
};

% ===== Main flow arrows =====
\draw[arrow] (cap.south) -- (cte.north);
\draw[arrow] (cte.south) |- (mdb.north);
\draw[arrow] (cte.south) |- (cdt.north);

\draw[arrow] (mdb.south) |- (spec.north);
\draw[arrow] (cdt.south) |- (proof.north);

\draw[softarrow] (spec.east) -- (impl.west);
\draw[softarrow] (impl.east) -- (proof.west);

\draw[softarrow] (resolve.north) -- (spec.south);
\draw[softarrow] (kernel.north) -- (proof.south);

\draw[framegray!45, dashed, line width=0.55pt] (gmgr.south) -- (boundary.north);

% ===== Side annotations =====
\node[note, left=10mm of gcap] (n1) {语义\\基底};
\node[note, left=10mm of gcte] (n2) {槽位\\抽象};
\node[note, left=10mm of gmdb] (n3) {图结构\\恢复};
\node[note, right=10mm of gcdt] (n4) {派生树\\保持};
\node[note, left=10mm of gmgr] (n5) {核心层\\主结论};
\node[note, right=10mm of gkernel] (n6) {接口层\\边界};

\draw[framegray!45, line width=0.5pt] (n1.east) -- (gcap.west);
\draw[framegray!45, line width=0.5pt] (n2.east) -- (gcte.west);
\draw[framegray!45, line width=0.5pt] (n3.east) -- (gmdb.west);
\draw[framegray!45, line width=0.5pt] (n4.west) -- (gcdt.east);
\draw[framegray!45, line width=0.5pt] (n5.east) -- (gmgr.west);
\draw[framegray!45, line width=0.5pt] (n6.west) -- (gkernel.east);

\end{tikzpicture}

\end{document}
